% eksperymentalne tłumaczenie na włoski

%

\section{Twierdzenie Menelausa} % lang-pl
\section{Teorema di Menelao} % lang-it

Znamy trzy twierdzenia o współliniowości: o prostej Eulera, o prostej Wallace'a-Simsona, o prostej Auberta. % lang-pl
Conosciamo tre teoremi di allineamento: sulla retta di Eulero, sulla retta di Wallace-Simson e sulla retta di Aubert. % lang-it

\begin{proposition}
[twierdzenie Menelaosa] % lang-pl
[teorema di Menelao] % lang-it
	...
	Wówczas punkty $K, L, M$ są współliniowe wtedy i tylko wtedy, gdy zachodzi % lang-pl
	Allora i punti $K$, $L$ e $M$ sono allineati se e solo se vale % lang-it
	\begin{equation}
		[AMB] [BKC] [CLA] = -1.
	\end{equation}
	\index{twierdzenie!Menelaosa} % lang-pl
	\index{teorema!di Menelao} % lang-it
\end{proposition}

% Eves s. 248

\index[persons]{Menelaos z Aleksandrii}% % lang-pl
\index[persons]{Menelao di Alessandria}% % lang-it

Brak pewności co do tego, kto będzie pierwszym odkrywcą tego. % lang-pl
Ale najstarsze ocalałe do czasów późniejszych wyjaśnienie pojawia się u Menelaosa (Σφαιρική) jako lemat, użyteczny w sferycznej wersji twierdzenia. % lang-pl
Non vi è certezza su chi ne sia stato il primo scopritore. % lang-it
Tuttavia, la più antica esposizione giunta fino a noi compare nella \emph{Sphaerica} (Σφαιρική) di Menelao come lemma utile per la versione sferica del teorema. % lang-it

Piszą o nim Guzicki \cite[s. 84]{guzicki_2021}; Hartshorne \cite[s. 180, 181]{hartshorne_2000}; Audin \cite[s. 38]{audin_2003}; Zetel \cite[s. 43]{zetel_2020}. % lang-pl
Uogólnieniem twierdzenia Menelaosa będzie twierdzenie Carnota (udowodnione w 1803 roku w \emph{,,Géométrie de position}''), że iloczyn długości odcinków wyznaczonych przez prostą na bokach wielokąta i niemających wspólnych końców jest równy iloczynowi długości pozostałych odcinków; % lang-pl
Ne scrivono Guzicki \cite[p. 84]{guzicki_2021}; Hartshorne \cite[pp. 180, 181]{hartshorne_2000}; Audin \cite[p. 38]{audin_2003} e Zetel \cite[p. 43]{zetel_2020}. % lang-it
Una generalizzazione del teorema di Menelao è il teorema di Carnot (dimostrato nel 1803 nella \emph{``Géométrie de position''}), secondo cui il prodotto delle lunghezze dei segmenti determinati da una retta sui lati di un poligono e privi di estremi comuni è uguale al prodotto delle lunghezze dei segmenti rimanenti; % lang-it
Pisze o nim Zetel w formie ćwiczenia \cite[s. 45]{zetel_2020}. % lang-pl
Ne scrive Zetel come esercizio \cite[p. 45]{zetel_2020}. % lang-it
% https://encyclopediaofmath.org/wiki/Carnot_theorem
% Zetel 45.

% TODO: https://en.wikipedia.org/wiki/Menelaus%27s_theorem
% It is uncertain who actually discovered the theorem; however, the oldest extant exposition appears in Spherics by Menelaus. In this book, the plane version of the theorem is used as a lemma to prove a spherical version of the theorem.

% TODO: https://www.deltami.edu.pl/2011/03/twierdzenie-menelaosa/

%